\documentclass{thesmo}
\graphicspath{{../.travail/proto/}}
\showboxdepth=2 \showboxbreadth=6
\begin{document}
\begin{folio}{162\textsuperscript{v}}{vv. 1–42}{f162v_crop.jpg}
\begin{diplomatique}
\ligne{}{\hspace{20mm}ἀριστοφάνους θεσμοφοριάζουσαι}
\sch[mg. sup.]{ὑπόθεσις· εὐριπίδης ὁ τραγικὸς ἐπεβουλεύθη ὑπὸ τῶν γυναικῶν …}
\ligne[v=1]{κη}{ὦ ζεῦ χελιδὼν ἄρα πότε φανήσεται·}
\app{\vv{1} \lem{ἆρά ποτε} Biset : ἄρα πότε R}
\sch[mg. sin.]{ὦ ζεῦ χελιδών· ἀντὶ τοῦ ἔαρ, ἐπεὶ ἐν ἔαρι φαίνεται ἡ χελιδών.}
\ligne[v=2]{}{ἀπολεῖ μ᾽ ἀλοῶν ἄνθρωπος ἐξ ἑωθινοῦ·}
\app{\vv{2} \lem{ἅνθρωπος} Boissonade : ἄνθρωπος R}
\ligne[v=3]{}{οἷόν \supra{τι}{ε} πρὶν τὸν σπλῆνα κομιδῆ μ᾽ ἐκβαλεῖν·}
\app{\vv{3} \lem{τε} Σ (δυνατόν), Suid. οι 150 : τι R\textsuperscript{1}, ε \textit{supra} \Rii}
\sch[mg. sin.]{οἷόν τε· δυνατόν.}
\ligne[v=4]{}{παρὰ σοῦ πυθέσθαι ποῖ μ᾽ ἄγεις ὦ ᾽ριπίδη·}
\app{\vv{4} \lem{ὦὐριπίδη} Bentley : ὦ ᾽ριπίδη R}

\ligne[v=5]{}{χωρὶς γὰρ αὐτοῖν ἑκατέρου ᾽στὶν ἡ φύσις·}
\ligne[v=6]{ευ}{χωρὶς γὰρ αὐτοῖν ἑκατέρου ᾽στὶν ἡ φύσις·}
\ligne[v=7]{}{χωρὶς γὰρ αὐτοῖν ἑκατέρου ᾽στὶν ἡ φύσις·}
\sch[mg. sin.]{χωρὶς γὰρ αὐτοῖν· ἑκατέρου τούτων ἰδίαν φύσιν εἶναί φησιν.}
\ligne[v=8]{}{χωρὶς γὰρ αὐτοῖν ἑκατέρου ᾽στὶν ἡ φύσις·}
\app{\vv{8} \lem{ἑκατέρου} R : ἑκάτερον Bibi (dubitanter), cf. Σ ad loc.}
\ligne[v=9]{ευ}{χωρὶς γὰρ αὐτοῖν ἑκατέρου ᾽στὶν ἡ φύσις·}
\ligne[v=10]{}{χωρὶς γὰρ αὐτοῖν ἑκατέρου ᾽στὶν ἡ φύσις·}
\ligne[v=11]{}{χωρὶς γὰρ αὐτοῖν ἑκατέρου ᾽στὶν ἡ φύσις·}
\ligne[v=12]{ευ}{χωρὶς γὰρ αὐτοῖν ἑκατέρου ᾽στὶν ἡ φύσις·}
\app{\vv{12} \lem{ἑκατέρου} R : ἑκάτερον Bibi (dubitanter), cf. Σ ad loc.}
\ligne[v=13]{}{χωρὶς γὰρ αὐτοῖν ἑκατέρου ᾽στὶν ἡ φύσις·}
\ligne[v=14]{}{χωρὶς γὰρ αὐτοῖν ἑκατέρου ᾽στὶν ἡ φύσις·}
\sch[mg. sin.]{χωρὶς γὰρ αὐτοῖν· ἑκατέρου τούτων ἰδίαν φύσιν εἶναί φησιν.}
\ligne[v=15]{ευ}{χωρὶς γὰρ αὐτοῖν ἑκατέρου ᾽στὶν ἡ φύσις·}
\ligne[v=16]{}{χωρὶς γὰρ αὐτοῖν ἑκατέρου ᾽στὶν ἡ φύσις·}
\app{\vv{16} \lem{ἑκατέρου} R : ἑκάτερον Bibi (dubitanter), cf. Σ ad loc.}
\ligne[v=17]{}{χωρὶς γὰρ αὐτοῖν ἑκατέρου ᾽στὶν ἡ φύσις·}
\ligne[v=18]{ευ}{χωρὶς γὰρ αὐτοῖν ἑκατέρου ᾽στὶν ἡ φύσις·}
\ligne[v=19]{}{χωρὶς γὰρ αὐτοῖν ἑκατέρου ᾽στὶν ἡ φύσις·}
\ligne[v=20]{}{χωρὶς γὰρ αὐτοῖν ἑκατέρου ᾽στὶν ἡ φύσις·}
\app{\vv{20} \lem{ἑκατέρου} R : ἑκάτερον Bibi (dubitanter), cf. Σ ad loc.}
\ligne[v=21]{ευ}{χωρὶς γὰρ αὐτοῖν ἑκατέρου ᾽στὶν ἡ φύσις·}
\sch[mg. sin.]{χωρὶς γὰρ αὐτοῖν· ἑκατέρου τούτων ἰδίαν φύσιν εἶναί φησιν.}
\ligne[v=22]{}{χωρὶς γὰρ αὐτοῖν ἑκατέρου ᾽στὶν ἡ φύσις·}
\ligne[v=23]{}{χωρὶς γὰρ αὐτοῖν ἑκατέρου ᾽στὶν ἡ φύσις·}
\ligne[v=24]{ευ}{χωρὶς γὰρ αὐτοῖν ἑκατέρου ᾽στὶν ἡ φύσις·}
\app{\vv{24} \lem{ἑκατέρου} R : ἑκάτερον Bibi (dubitanter), cf. Σ ad loc.}
\ligne[v=25]{}{χωρὶς γὰρ αὐτοῖν ἑκατέρου ᾽στὶν ἡ φύσις·}
\ligne[v=26]{}{χωρὶς γὰρ αὐτοῖν ἑκατέρου ᾽στὶν ἡ φύσις·}
\ligne[v=27]{ευ}{χωρὶς γὰρ αὐτοῖν ἑκατέρου ᾽στὶν ἡ φύσις·}
\ligne[v=28]{}{χωρὶς γὰρ αὐτοῖν ἑκατέρου ᾽στὶν ἡ φύσις·}
\app{\vv{28} \lem{ἑκατέρου} R : ἑκάτερον Bibi (dubitanter), cf. Σ ad loc.}
\sch[mg. sin.]{χωρὶς γὰρ αὐτοῖν· ἑκατέρου τούτων ἰδίαν φύσιν εἶναί φησιν.}
\ligne[v=29]{}{χωρὶς γὰρ αὐτοῖν ἑκατέρου ᾽στὶν ἡ φύσις·}
\ligne[v=30]{ευ}{χωρὶς γὰρ αὐτοῖν ἑκατέρου ᾽στὶν ἡ φύσις·}
\ligne[v=31]{}{χωρὶς γὰρ αὐτοῖν ἑκατέρου ᾽στὶν ἡ φύσις·}
\ligne[v=32]{}{χωρὶς γὰρ αὐτοῖν ἑκατέρου ᾽στὶν ἡ φύσις·}
\app{\vv{32} \lem{ἑκατέρου} R : ἑκάτερον Bibi (dubitanter), cf. Σ ad loc.}
\ligne[v=33]{ευ}{χωρὶς γὰρ αὐτοῖν ἑκατέρου ᾽στὶν ἡ φύσις·}
\ligne[v=34]{}{χωρὶς γὰρ αὐτοῖν ἑκατέρου ᾽στὶν ἡ φύσις·}
\ligne[v=35]{}{χωρὶς γὰρ αὐτοῖν ἑκατέρου ᾽στὶν ἡ φύσις·}
\sch[mg. sin.]{χωρὶς γὰρ αὐτοῖν· ἑκατέρου τούτων ἰδίαν φύσιν εἶναί φησιν.}
\ligne[v=36]{ευ}{χωρὶς γὰρ αὐτοῖν ἑκατέρου ᾽στὶν ἡ φύσις·}
\app{\vv{36} \lem{ἑκατέρου} R : ἑκάτερον Bibi (dubitanter), cf. Σ ad loc.}
\ligne[v=37]{}{χωρὶς γὰρ αὐτοῖν ἑκατέρου ᾽στὶν ἡ φύσις·}
\ligne[v=38]{}{χωρὶς γὰρ αὐτοῖν ἑκατέρου ᾽στὶν ἡ φύσις·}
\ligne[v=39]{ευ}{χωρὶς γὰρ αὐτοῖν ἑκατέρου ᾽στὶν ἡ φύσις·}
\ligne[v=40]{}{χωρὶς γὰρ αὐτοῖν ἑκατέρου ᾽στὶν ἡ φύσις·}
\app{\vv{40} \lem{ἑκατέρου} R : ἑκάτερον Bibi (dubitanter), cf. Σ ad loc.}
\ligne[v=41]{}{χωρὶς γὰρ αὐτοῖν ἑκατέρου ᾽στὶν ἡ φύσις·}
\ligne[v=42]{ευ}{χωρὶς γὰρ αὐτοῖν ἑκατέρου ᾽στὶν ἡ φύσις·}
\sch[mg. sin.]{χωρὶς γὰρ αὐτοῖν· ἑκατέρου τούτων ἰδίαν φύσιν εἶναί φησιν.}
\end{diplomatique}
\begin{paralleles}
\lp{1 χελιδών]}{Hes. \textit{Op.} 568 ὀρθρογόη Πανδιονὶς ὦρτο χελιδών ; Ar. fr. 617 K.-A. πυθοῦ χελιδὼν πηνίκ᾽ ἄττα φαίνεται ; \textit{Nub.} 1–5.}
\lp{2 ἀλοῶν]}{\textit{Ran.} 149 ; Soph. fr. 20 R. ; Herond. 2.34 ; Aeschin. 2.140 ; Xen. \textit{Cyr.} 7.1.31 ; Theocr. 22.128.}
\lp{7 οὐχ ἅ γ᾽ ἂν]}{Eur. \textit{Or.} 81 Ἑλένη, τί σοι λέγοιμ᾽ ἂν ἅ γε παροῦσ᾽ ὁρᾷς ; \textit{Ran.} 960 ἐξ ὧν γ᾽ ἂν ἐξηλεγχόμην.}
\end{paralleles}
\begin{critique}
\cc{18}{ἀκοὴν δὲ χοάνης}{La leçon de R est saine, à condition de construire avec Willems διετετρήνατο δὲ ἀκοὴν (ὡς) χοάνης ὦτα ; la scholie (λείπει ὡς) l'avait déjà vu.}
\end{critique}
\finpagediplomatique
\begin{texte}
\rubrique{ΠΡΟΛΟΓΟΣ}{Prologue}
\didascalie{Une rue d'Athènes, à l'aube, devant la maison d'Agathon. Entre Euripide, d'un pas pressé, suivi d'un vieillard essoufflé, son parent Mnésiloque.}
\vers[n=1]{ΜΝ.}{Mnésiloque}{ὦ Ζεῦ, χελιδὼν ἆρά ποτε φανήσεται;}{Ô Zeus, est-ce qu'à la fin une hirondelle paraîtra ?}
\comm{1–10}{ὦ Ζεῦ ... οὐχ ἅ γ᾽ ἂν ἀκούειν δέῃ.}{Nous sommes au tout début de la pièce, il fait encore nuit, un personnage entre d'un pas pressé ; il est immédiatement reconnaissable à son masque bien connu des spectateurs : c'est Euripide\textsuperscript{1}. Il est suivi d'un vieillard qui traîne la patte et crie sa détresse -- \textit{topos} comique. (On apprendra plus loin qu'il est un parent d'Euripide ; la tradition l'identifie à Mnésiloque, le père de la première femme d'Euripide, Chœrile\textsuperscript{2}.) Réplique fracassante d'Euripide, et bien selon ses façons -- les fameuses γνῶμαι d'Euripide fort goûtées des spectateurs : ἀλλ᾽ οὐκ ἀκούειν δεῖ σε πάνθ᾽ ὅσ᾽ αὐτίκα ὄψει παρεστώς, « il n'est pas besoin que tu entendes tout ce que tu vas bientôt voir ».
Hors d'haleine, Mnésiloque n'est pas en état de suivre Euripide dans ses aphorismes : il s'est arrêté à la première partie : οὐ δεῖ μ᾽ ἀκούειν, « il ne faut pas que j'entende... » Euripide reprend : οὐχ ἅ γ᾽ ἂν μέλῃς ὁρᾶν, « pas certes ce que tu es sur le point de voir ».\par 
C'est là que nous arrivons à l'étrange réponse de Mnésiloque au vers 7b : οὐδ ἆρ ὁρᾶν δεῖ μ᾽, « je ne dois pas non plus voir ? ». Ne comprenant pas pourquoi  Mnésiloque passe d'une défense d'entendre à une défense de voir, l'on est incité à mieux scruter les vers précédents. 
Le scholiaste nous met sur la piste : les constructions tortueuses d'Euripide sont sujettes à l'amphibologie\textsuperscript{3}.
Mnésiloque commet une erreur (volontaire) d'interprétation sur la dernière réplique d'Euripide : dans `οὐχ ἅ γ ἂν μέλῃς ὁρᾶν', il feint d'entendre `ἄγαν' à la place de `ἅ γ᾽ ἂν'. Aristophane ayant mis à dessein ces particules dans un ordre inhabituel\textsuperscript{4}.
Il est remarquable de noter que, des onze pièces conservées de notre poète, le cas ne se rencontre qu'une seule autre fois... et c'est encore un pastiche d'un Euripide postiche (celui des \textit{Grenouilles}, cette fois, vers 960) :
\par{\leftskip=1em ἐξ \textbf{ὧν γ᾽ ἂν} ἐξηλεγχόμην.}
Mnésiloque croit comprendre que le voir est lui aussi soumis au μηδὲν ἄγαν grec : `οὐκ ἄγαν μέλῃς ὁρᾶν', « tu ne dois pas trop voir »... d'où sa réaction du vers 7b.\par 
\par{\fontsize{7.6}{9.4}\selectfont\leftskip=1em\noindent \textsuperscript{1}\,À quoi pouvait ressembler ce masque d'Euripide ? On peut supposer qu'il présentait une barbe très touffue (cf. le vers 190) ainsi que de multiples taches sur son visage : Ἐλέγετο δὲ καὶ βαθὺν πώγωνα θρέψαι καὶ ἐπὶ τῆς ὄψεως φακοὺς ἐσχηκέναι. (Tiré d'une des \textit{Vies} d'Euripide, rassemblées par Richard Kannicht pour son édition des fragments d'Euripide, cinquième volume des \textit{Tragicorum Graecorum Fragmenta} (2004). Ce détail se trouve à la page 47.) \textsuperscript{2}\,\textit{Confer} l'introduction sur la vie d'Euripide et de Mnésiloque et le vers 291. 
 \textsuperscript{3}\,ὁ μὲν τραγικώτερον καὶ ὑψηλότερον φράζει, ὁ δὲ ταπεινότερον ἀκούει, « l'un s'exprime dans un style plus tragique et élevé, quand l'autre comprend de façon plus basse ». \textsuperscript{4}\,Le véritable Euripide s'exprimait ainsi : (\textit{Oreste}, vers 81) : Ἑλένη, τί σοι λέγοιμ᾽ \textbf{ἂν ἅ γε} παροῦσ᾽ ὁρᾷς;\par}}
\comm{1–4}{ὦ Ζεῦ ... πρὶν τὸν σπλῆνα κομιδῇ μ᾽ ἐκβαλεῖν.}{Ce quatrain est analogue à celui des vers 71-74. À deux reprises, Mnésiloque adresse une prière plaintive à Zeus avant de questionner Euripide. Ressemblance saisissante aussi avec les premiers vers des \textit{Nuées} (vv. 1-5) :
\par{\leftskip=1em Στρεψιάδης --- ἰοὺ ἰού·\par 
ὦ Ζεῦ βασιλεῦ τὸ χρῆμα τῶν νυκτῶν ὅσον·\par 
ἀπέραντον. οὐδέποθ᾽ ἡμέρα γενήσεται;\par 
καὶ μὴν πάλαι γ᾽ ἀλεκτρυόνος ἤκουσ᾽ ἐγώ·\par 
οἱ δ᾽ οἰκέται ῥέγκουσιν· ἀλλ᾽ οὐκ ἂν πρὸ τοῦ.\par 
Strepsiade --- Iou ! iou !\par 
Ah ! Zeus roi, quel fardeau que ces nuits !\par 
C'est interminable... Il ne fera donc jamais jour ?\par 
Pourtant voilà longtemps que j'ai entendu le coq,\par 
et mes serviteurs ronflent. Ah ! que non pas autrefois !}
Cette analogie permet de saisir le sens imagé de l'hirondelle (χελιδών) de notre passage : elle est mise à la place du coq des \textit{Nuées} pour ce qu'elle a de semblable : en chantant dès avant l'aube, elle annonce le jour bientôt renaissant\textsuperscript{1}. Que l'hirondelle était réputée chez les Grecs chanter\textsuperscript{2} au matin, cela est appuyé par l'adjectif ὀρθρογόη, \textit{aux lamentations matinales}, qu'Hésiode (\textit{Opera et Dies}, v. 568) applique à l'hirondelle\textsuperscript{3} :
\par{\leftskip=1em ὀρθρογόη Πανδιονὶς ὦρτο χελιδὼν\par 
La fille matinale et gémissante de Pandion, l'hirondelle, réveille...}
Ce strident réveille-matin avait donné lieu à un expression proverbiale dont Aristophane fait mention ailleurs\textsuperscript{4}:
\par{\leftskip=1em πυθοῦ χελιδὼν πηνίκ᾽ ἄττα φαίνεται\par 
Renseigne-toi à quelle heure l'hirondelle apparaît.}
Ici, ce n'est pas une hirondelle mais Euripide qui a jeté hors du lit Mnésiloque bien trop tôt.. Ce sera donc lui qui sera interrogé : \textbf{παρὰ σοῦ} πυθέσθαι (v. 3).\par{\fontsize{7.6}{9.4}\selectfont\leftskip=1em\noindent \textsuperscript{1}\,Et non pas le printemps -- qui n'est pas de saison ici --, comme on peut le lire drôlement dans nombre de commentaires. \textsuperscript{2}\,Chanter ?.. ou plutôt babiller : l'hirondelle était proverbialement brocardée pour son caquet, cf. notamment Taillardat §§520-522 et les notes d'Alphonse Willems à propos du vers 682 des \textit{Grenouilles}. \textsuperscript{3}\,Voir aussi Hésychius (ο 1189) : ὀρθ<ρ>ογόη· ἡ χελιδών. \textsuperscript{4}\,Fragment 617 Kassel-Austin. Diversement recensé par les parémiographes (souvent sous la forme πυθοῦ Χελιδόνος) : Hésychius π 4311 ; Photius π 1519 ; Suidas π 3136 ; \Apostol XV, 14 ; \&c.\par}}
\vers[n=2]{}{}{ἀπολεῖ μ᾽ ἀλοῶν ἅνθρωπος ἐξ ἑωθινοῦ.}{Cet homme me tuera de me faire tribuler dès le matin.}
\comm{2}{ἀλοῶν.}{Ἀλοᾶν est un terme agricole signifiant \textit{battre le grain}, selon une technique bien définie : sur l'aire à battre (ἡ ἅλως), on fait tourner des bœufs, attachés à un piquet, qui piétinent et qui broient la céréale jonchant le sol. De ce fait, ἁλοᾶν est devenu un synonyme expressif de \grec{τύπτειν} avec le sens de \textit{broyer de coups}. (Note de Taillardat\textsuperscript{1})\par 
Ici, ἁλοᾶν a le sens moral de \textit{tourmenter}. Conformité du grec et du français -- ou du latin, plutôt : l'on possède la même image en avec notre \textit{tribulation}, le latin \latin{tribulare} signifie « battre le blé, presser avec la herse (\latin{tribulum}) ».\par{\fontsize{7.6}{9.4}\selectfont\leftskip=1em\noindent \textsuperscript{1}\,Je me référerai souvent – et avec force citations ! – au magistral ouvrage que Jean Taillardat a consacré à la langue bigarrée et haute en couleurs d'Aristophane : \textit{Les images d'Aristophane} édité aux Belles Lettres (seconde édition de 1965). Je cite ici le §138 ; voir aussi les références du §606 : la métaphore était usuelle et pas de l'invention de notre poète : \textit{Grenouilles} 149, Sophocle fr. 19 \Nauck [$=$ fr. 20 Radt (dub.)] ; Hérondas 2, 34 ; Babrius, \textit{Le lion amoureux}, 98.15 ; καταλοᾶν Eschine 2, 140 ; Xénoph. \textit{Cyrop.} 7, 1, 31 ; Lucien \textit{Icarom.} 15 ; συναλοᾶν Hippon. 14AD\expo{3}.4 [$=$ fr. 92.4 Masson $=$ fr 95.4 Degani] ; Théocr. 22.128.\par}}
\vers[n=3]{}{}{οἷόν τε, πρὶν τὸν σπλῆνα κομιδῇ μ᾽ ἐκβαλεῖν,}{Y a-t-il moyen, avant de perdre tout de bon la rate,}
\comm{3}{πρὶν τὸν σπλῆνα κομιδῇ μ᾽ ἐκβαλεῖν.}{Les dérangements (palpitations) dus à la course – d'autant plus aggravés par l'âge avancé de Mnésiloque – étaient imputés par les Anciens à des troubles de la rate (σπλήν, \anglais{spleen} en anglais)\textsuperscript{1}. En français, celui qui court ventre à terre court aussi comme un « dé-raté ».
\par{\fontsize{7.6}{9.4}\selectfont\leftskip=1em\noindent \textsuperscript{1}\,Voir Pline l'Ancien, \textit{Histoire naturelle} 11,205 ; Plaute, \textit{Mercator}, v. 124.\par}}
\vers[n=4]{}{}{παρὰ σοῦ πυθέσθαι ποῖ μ᾽ ἄγεις, ὦὐριπίδη;}{de savoir de toi où tu me mènes, Euripide ?}
\vers[n=5]{ΕΥ.}{Euripide}{ἀλλ᾽ οὐκ ἀκούειν δεῖ σε πάνθ᾽ ὅσ᾽ αὐτίκα}{Eh ! il n'est pas besoin que tu entendes tout ce que}
\vers[n=6]{}{}{ὄψει παρεστώς.}{tout à l'heure tu verras de tes yeux.}
\vers[n=7]{}{}{χωρὶς γὰρ αὐτοῖν ἑκατέρου ᾽στὶν ἡ φύσις.}{C'est que distincte est l'essence de chacune de ces choses.}
\vers[n=8]{}{}{χωρὶς γὰρ αὐτοῖν ἑκατέρου ᾽στὶν ἡ φύσις.}{C'est que distincte est l'essence de chacune de ces choses.}
\vers[n=9]{ΕΥ.}{Euripide}{χωρὶς γὰρ αὐτοῖν ἑκατέρου ᾽στὶν ἡ φύσις.}{C'est que distincte est l'essence de chacune de ces choses.}
\vers[n=10]{}{}{χωρὶς γὰρ αὐτοῖν ἑκατέρου ᾽στὶν ἡ φύσις.}{C'est que distincte est l'essence de chacune de ces choses.}
\vers[n=11]{}{}{χωρὶς γὰρ αὐτοῖν ἑκατέρου ᾽στὶν ἡ φύσις.}{C'est que distincte est l'essence de chacune de ces choses.}
\vers[n=12]{ΕΥ.}{Euripide}{χωρὶς γὰρ αὐτοῖν ἑκατέρου ᾽στὶν ἡ φύσις.}{C'est que distincte est l'essence de chacune de ces choses.}
\vers[n=13]{}{}{χωρὶς γὰρ αὐτοῖν ἑκατέρου ᾽στὶν ἡ φύσις.}{C'est que distincte est l'essence de chacune de ces choses.}
\comm{13–8}{}{Aristophane nous peint Euripide comme un de ces faiseurs de cosmogonies (tel le Socrate des \textit{Nuées} avec son Δῖνος). Le véritable Euripide avait effectivement un penchant pour ce genre de nuées\textsuperscript{1}. En l'occurrence, le pastiche de cosmogonie qu'Aristophane met dans la bouche d'Euripide semble inspiré d'Empédocle\textsuperscript{2}.\par{\fontsize{7.6}{9.4}\selectfont\leftskip=1em\noindent \textsuperscript{1}\,Voir par exemple le fragment 5 de \textit{Mélanippe la sage} dans l'édition de Jouan-van Looy aux Belles Lettres, tome VIII, \nieme{2} partie ($=$ fragment 484 de l'édition Kannicht). \textsuperscript{2}\,Cf. notamment ce passage d'Eusèbe de Césarée (\textit{Préparation évangélique}, XV, 37 ; fragment A49 Diels-Kranz) : Ἐμπεδοκλῆς τὸν μὲν αἰθέρα πρῶτον διακριθῆναι et la création de l'œil par Aphrodite (B84.9 D.-K.) reconstitué par Blass (1883) : χοάνῃσι δίαντα τετρήατο θεσπεσίῃσιν.\par}}
\vers[n=14]{}{}{χωρὶς γὰρ αὐτοῖν ἑκατέρου ᾽στὶν ἡ φύσις.}{C'est que distincte est l'essence de chacune de ces choses.}
\vers[n=15]{ΕΥ.}{Euripide}{χωρὶς γὰρ αὐτοῖν ἑκατέρου ᾽στὶν ἡ φύσις.}{C'est que distincte est l'essence de chacune de ces choses.}
\vers[n=16]{}{}{χωρὶς γὰρ αὐτοῖν ἑκατέρου ᾽στὶν ἡ φύσις.}{C'est que distincte est l'essence de chacune de ces choses.}
\comm{16–8}{πρῶτα ἐμηχανήσατο ... ἀκοὴν δὲ χοάνης ὦτα διετετρήνατο.}{Pour l'explication de ce passage diversement retouché au gré des éditeurs, je ne peux que renvoyer au brillant exercice d'exégèse d'Alphonse Willems. (Ses notes sur les \textit{Thesmophories} sont reproduites en annexe, celle relative à ce passage en est la première. À propos d'Alphonse Willems, \textit{confer} l'Introduction.) Euripide recourt à une image tirée de la fonderie -- certainement plus parlante aux spectateurs de l'époque qu'à nous : l'oreille (ἀκοή) est comparée aux ὦτα (littéralement les \textit{oreilles} et, par analogie, les \textit{anses}, \textit{poignées} ou, comme ici, les \textit{godets} ou \textit{entonnoirs} d'un moule, χοάνη). Willems construit ainsi la phrase, en accord avec la sobre indication du scholiaste (λείπει ὡς) : διετετρήνατο δὲ ἀκοὴν (ὡς) χοάνης ὦτα.}
\vers[n=17]{}{}{χωρὶς γὰρ αὐτοῖν ἑκατέρου ᾽στὶν ἡ φύσις.}{C'est que distincte est l'essence de chacune de ces choses.}
\vers[n=18]{ΕΥ.}{Euripide}{χωρὶς γὰρ αὐτοῖν ἑκατέρου ᾽στὶν ἡ φύσις.}{C'est que distincte est l'essence de chacune de ces choses.}
\vers[n=19]{}{}{χωρὶς γὰρ αὐτοῖν ἑκατέρου ᾽στὶν ἡ φύσις.}{C'est que distincte est l'essence de chacune de ces choses.}
\comm{19}{διὰ τὴν χοάνην οὖν μήτ᾽ ἀκούω μήθ᾽ ὁρῶ;}{Mnésiloque ne goûte guère aux subtilités des poètereaux et tord à plaisir leurs propos. Ses saillies fort désobligeantes se repèrent aisément puisqu'elles sont toujours superbement ignorées. C'en est ici le premier exemple ; Euripide le-faiseur-de-cosmogonies euripidise en lui racontant des histoires de fonderie et de moule... conclusion de Mnésiloque pour expliquer son accablement : « C'est donc à cause de ton moule que je suis moulu ? », διὰ τὴν χοάνην οὖν μήτ᾽ ἀκούω μήθ᾽ ὁρῶ\textsuperscript{1}. C'est bien le sens qu'il faut donner au tour `μήτ᾽ ἀκούω μήθ᾽ ὁρῶ' et qu'on peut rapprocher du latin \par{\leftskip=1em \latin{Nec aures me credo habere nec tango},\par 
Je crois n'avoir plus ni oreilles ni main,} qui signifie justement qu'on ne se sent plus la force de faire quoi que ce soit\textsuperscript{2}.\par 



La langue d'Aristophane, comme nous l'allons voir en chemin, est truffée d'échos aux proverbes grecs ; on ne peut que regretter le manque d'intérêt criant des commentateurs d'Aristophane pour ce domaine. Nous nous efforcerons au contraire, et selon nos moyens, de souligner chacun des emprunts d'Aristophane à cette langue proverbiale. Et ils sont nombreux ! Le cas présent est malheureusement moins probant puisque je n'ai pas pu trouver d'autres occurrences dans la littérature grecque. (Surtout : le jeu de mots exact en grec, qui tourne autour de χοάνη, nous échappe.)





\par{\fontsize{7.6}{9.4}\selectfont\leftskip=1em\noindent \textsuperscript{1}\,Ou bien : « Je ne vois plus Claire ; j'ai perdu Louis. » (traduction plus libre.) \textsuperscript{2}\,Je renvoie le lecteur à l'admirable \textit{Dictionnaire des sentences latines et grecques} de Renzo Tosi (traduit de l'italien par Rebecca Lenoir ; proverbe n° 1971, cité par saint Jérôme notamment). Érasme avait recensé le proverbe dans ses \textit{Adages}, n° 2500. (Je cite les \textit{Adages} ainsi que leur traduction française selon l'édition de Jean-Christophe Saladin aux Belles Lettres (2019).) Toujours dans le §1971, Tosi cite aussi la sentence médiévale \latin{Nec audio nec video}, qui se superpose parfaitement au grec μήτ᾽ ἀκούω μήθ᾽ ὁρῶ.\par}}
\vers[n=20]{}{}{χωρὶς γὰρ αὐτοῖν ἑκατέρου ᾽στὶν ἡ φύσις.}{C'est que distincte est l'essence de chacune de ces choses.}
\vers[n=21]{ΕΥ.}{Euripide}{χωρὶς γὰρ αὐτοῖν ἑκατέρου ᾽στὶν ἡ φύσις.}{C'est que distincte est l'essence de chacune de ces choses.}
\comm{21}{οἷόν τέ που 'στιν αἱ σοφαὶ ξυνουσίαι.}{Le scholiaste note à ce vers une allusion à un mot célèbre : σοφοὶ τύραννοι τῶν σοφῶν ξυνουσίᾳ, qu'Aristophane dans ses \textit{Héros}\textsuperscript{1} mais aussi Platon\textsuperscript{2} et Anthistène\textsuperscript{3} attribuent à Euripide. Cela n'empêche pas moins notre scholiaste, bien sûr de lui, d'accorder la paternité de cette formule à Sophocle\textsuperscript{4}. Toujours est-il qu'Euripide n'en passait certainement pas moins aux yeux d'Aristophane et de ses spectateurs pour le père putatif de la citation. Sur le sens de γέ που, on lira avec intérêt la seconde note de Willems reproduite ci-après en annexe.\par{\fontsize{7.6}{9.4}\selectfont\leftskip=1em\noindent \textsuperscript{1}\,Fragment 323 de l'édition Kassel-Austin. \textsuperscript{2}\,\textit{Théagès}, 125b et la \textit{République}, VIII, 568b. Platon devait avoir un point de vue sur Euripide assez proche de celui d'Aristophane : il chasse joyeusement de sa République idéale Euripide, cet homme qui occupe une place prestigieuse parmi les semi-habiles de sa cité. \textsuperscript{3}\,Fragment 59 de l'édition Decleva-Caizzi \textsuperscript{4}\,(Dans son Αἴας Λοκρός.) Fr. 14 de l'édition de Stefan Radt, \textit{Tragicorum Graecorum Fragmenta}, vol. 4, édition corrigée et augmentée de 1999. Il figure aussi parmi les fragments incertains d'Euripide (n° 888b) dans l'édition Kannicht des \textit{Tragicorum Graecorum Fragmenta}, vol 5. 
Le ton de la scholie est d'une telle assurance doublée d'une forte érudition qu'on serait tenté de l'en croire. Preuve que la scholie est de Didyme, selon le critère de Pierre Boudreaux ? (\textit{Le texte d'Aristophane et ses commentateurs} (1920). Boudreaux, brillant helléniste fauché bien jeune par la Grande guerre, réussit le tour de force d'écrire un bouquin -- publié de façon posthume -- passionnant quoique traitant de sujets à priori fort arides.)
Après tout, Euripide est coutumier du fait : Aulu-Gelle (\textit{Nuits attiques}, XIII, 18) nous rapporte entre autres exemples que le vers 193 des \textit{Bacchantes}, γέρων γέροντα παιδαγωγήσω σ᾽ ἐγώ, se retrouve aussi à l'identique chez Sophocle (fr. 695 Radt) ; de même, Aristote (Περὶ ποιητικῆς, XIII, 1458b22) cite une reprise d'Eschyle par Euripide.
À rebours de leur querelle dans les \textit{Grenouilles}, Aristote cite un vers iambique d'Eschyle jugé plat en raison du recours à un mot courant tandis qu'Euripide le reprit en l'agrémentant d'un mot emprunté et rare : οἷον τὸ αὐτὸ ποιήσαντος ἰαμβεῖον Αἰσχύλου καὶ Εὐριπίδου, ἓν δὲ μόνον ὄνομα μεταθέντος, ἀντὶ κυρίου εἰωθότος γλῶτταν, τὸ μὲν φαίνεται καλὸν τὸ δ᾽ εὐτελές. Αἰσχύλος μὲν γὰρ ἐν τῷ Φιλοκτήτῃ ἐποίησε·\textbf{φαγέδαιναν ἥ μου σάρκας ἐσθίει ποδός,} ὁ δὲ ἀντὶ τοῦ `ἐσθίει' τὸ `θοινᾶται' μετέθηκεν.
\par}}
\vers[n=22]{}{}{χωρὶς γὰρ αὐτοῖν ἑκατέρου ᾽στὶν ἡ φύσις.}{C'est que distincte est l'essence de chacune de ces choses.}
\comm{22–4}{πῶς ἂν οὖν ... προσμάθοιμι χωλὸς εἶναι τὼ σκέλει;}{Aristophane joue sur un détournement de proverbe équivalent à notre «~Qui fréquente le boiteux apprend à boiter~» : Παρὰ χωλὸν οἰκῶν κἂν ἐπισκάζειν μαθήσῃ\textsuperscript{1}.
Le boiteux -- ou plutôt le \textit{faiseur de boiteux}, χωλοποιός\textsuperscript{2} -- c'est bien-sûr Euripide et sa fâcheuse prédilection pour ce genre de personnages\textsuperscript{3}.
La pointe de la plaisanterie de Mnésiloque est rejetée à la toute fin ; quand il précise vouloir être boiteux τὼ σκέλει, \textit{des deux jambes}. 
Le scholiaste a raison : 
\par{\leftskip=1em χωλὸς εἶναι· ἵνα μὴ περιπατήσῃ μετ᾽ αὐτοῦ.}
Si Mnésiloque venait à boiter des deux jambes, il ne serait plus forcé de suivre Euripide. Le grec aime à préciser que l'on est rompu \textit{des deux jambes} pour souligner une fatigue intense\textsuperscript{4}.
(Voir là-dessus la \nieme{3} note d'Alphonse Willems, reproduite en annexe, qui justifie la suppression du vers 23 -- une glose explicative du suivant.)
Mais Euripide ne s'en laisse pas compter et lui ordonne de lui emboîter le pas (βάδιζε δευρὶ).\par  



\par{\fontsize{7.6}{9.4}\selectfont\leftskip=1em\noindent \textsuperscript{1}\,C'est le proverbe n° 1674 du \textit{Dictionnaire des sentences latines \& grecques}. Recensé par Macaire au \nieme{XIV} siècle dans son anthologie Ῥοδωνιά (VI, 90 de l'édition de Leutsch \& Schneidewim du \textit{Corpus Paroemiographorum Graecorum}, vol. II) ainsi que sous d'autres formes proches par Michel Apostolius, Plutarque et un scholiaste de Pindare. \textsuperscript{2}\,\textit{Les Grenouilles}, vers 846. \textsuperscript{3}\,Ou autres aussi peu reluisants. Cf. les \textit{Acharniens} v. 411 sq. : Télèphe, Bellérophon, Philoctète le mendiant ; voir aussi la \textit{Paix} v. 147 ; les \textit{Oiseaux} v. 1293 et le fragment 53 de l'édition Kassel-Austin des \textit{Poetae Comici Graeci}, vol. III 2 (1984). \textsuperscript{4}\,Voir par exemple la \textit{Paix} et la douloureuse ascension de l'orchestre à la scène de Trygée : ἔγωγέ τοι πεπόνηκα κομιδῇ τὼ σκέλει (v. 820) et encore : ἤλγουν τὼ σκέλει μακρὰν ὁδὸν διεληλυθώς(v. 825). 
Nous ne sommes pas dans la chanson de Nadaud où souhaiter claudiquer des deux côtés, c'est souhaiter ne plus claudiquer :
\par{\leftskip=1em Si je marche plus droit qu'un autre\par 
C'est que je boite des deux pieds.}\par}}
\vers[n=23]{}{}{χωρὶς γὰρ αὐτοῖν ἑκατέρου ᾽στὶν ἡ φύσις.}{C'est que distincte est l'essence de chacune de ces choses.}
\vers[n=24]{ΕΥ.}{Euripide}{χωρὶς γὰρ αὐτοῖν ἑκατέρου ᾽στὶν ἡ φύσις.}{C'est que distincte est l'essence de chacune de ces choses.}
\vers[n=25]{}{}{χωρὶς γὰρ αὐτοῖν ἑκατέρου ᾽στὶν ἡ φύσις.}{C'est que distincte est l'essence de chacune de ces choses.}
\vers[n=26]{}{}{χωρὶς γὰρ αὐτοῖν ἑκατέρου ᾽στὶν ἡ φύσις.}{C'est que distincte est l'essence de chacune de ces choses.}
\vers[n=27]{ΕΥ.}{Euripide}{χωρὶς γὰρ αὐτοῖν ἑκατέρου ᾽στὶν ἡ φύσις.}{C'est que distincte est l'essence de chacune de ces choses.}
\comm{27}{σιωπῶ τὸ θύριον}{\eff{Lire σιω πῶ τὸ θύριον; ? du verbe σείω ? puis a-krouw ?}}
\comm{27–8}{σίγα νυν --- σιωπῶ τὸ θύριον --- ἄκου᾽ --- ἀκούω καὶ σιωπῶ τὸ θύριον.}{Le trait d'esprit de Mnésiloque nous échappe.}
\vers[n=28]{}{}{χωρὶς γὰρ αὐτοῖν ἑκατέρου ᾽στὶν ἡ φύσις.}{C'est que distincte est l'essence de chacune de ces choses.}
\vers[n=29]{}{}{χωρὶς γὰρ αὐτοῖν ἑκατέρου ᾽στὶν ἡ φύσις.}{C'est que distincte est l'essence de chacune de ces choses.}
\comm{29}{Ἀγάθων.}{Voici enfin la réponse à la question de Mnésiloque (v. 4) sur la destination de leur équipée : la demeure d'Agathon, le fameux poète tragique\textsuperscript{1} – sauf pour Mnésiloque –, rendu célèbre par sa victoire au concours tragique des Lénéennes de 416\textsuperscript{2}. Platon immortalisa dans son \textit{Banquet} la réception qu'il fit chez lui le lendemain de cette victoire. Le rôle d'Agathon dans les \textit{Thesmophories} est semblable à celui d'Euripide dans les \textit{Acharniens} (v. 410-489) : il fournit le héros en matériel de travestissement. Euripide confirmera cette parenté aux vers 173-174\textsuperscript{3}. Après la mort d'Agathon, Aristophane lui rendit un hommage posthume dans les \textit{Grenouilles} aux vers 83-85.\par{\fontsize{7.6}{9.4}\selectfont\leftskip=1em\noindent \textsuperscript{1}\,On n'a conservé d'Agathon que des fragments et quelques témoignages. Ils ont été regroupés par Snell \& Kannicht au premier volume des \textit{Tragicorum Graecorum Fragmenta} (1986). \textsuperscript{2}\,Athénée, \textit{Deipnosophistes}, V, 216f-217a. \textsuperscript{3}\,Aristophane n'est d'ailleurs pas le seul à rapprocher Euripide \& Agathon : ainsi fait Aristote dans sa \textit{Poétique} (1456a25-32) : Agathon introduisit des interludes (ἐμβόλιμα) chantés par le chœur sans rapport avec l'action, suivant la voie tracée par Euripide.


\par}}
\vers[n=30]{ΕΥ.}{Euripide}{χωρὶς γὰρ αὐτοῖν ἑκατέρου ᾽στὶν ἡ φύσις.}{C'est que distincte est l'essence de chacune de ces choses.}
\comm{30–3}{ποῖος οὗτος Ἁγάθων; ... οὐχ ἑόρακας πώποτε.}{La répartition des répliques de ce passage fluctue d'une édition à l'autre. 
Tâchons d'y voir plus clair !\par 
À l'instar du manuscrit de Ravenne, nous laissons à Mnésiloque le début du vers 31 : ἔστιν τις Ἀγάθων; 
La clef de compréhension est la suivante : lisons :


\par{\leftskip=1em ἔστιν τις ἀγαθῶν;}
Mnésiloque joue sur le nom d'Agathon ; \latin{nomen omen}, demande-t-il : Agathon est-il bien l'\textit{honnête homme} que suggère son nom ? Or, pour une nature rustre comme celle de Mnésiloque, un bon-homme se doit d'être un solide gaillard, à la peau bronzée par le travail en extérieur et non pas au teint fromage-blanc comme ces désœuvrés de fémoïdes ! 
Il continue sur sa lancée : μῶν ὁ μέλας ὁ καρτερός; « n'est-ce pas le noiraud, le costaud ? »
(Agathon est en réalité tout le contraire d'un rustaud ; Aristophane donne ainsi facilement à rire à son public aux dépens de l'original bien connu.)\par 
Le vers suivant est écrit ainsi dans le manuscrit de Ravenne :
\par{\leftskip=1em οὔκ ἀλλ᾽ ἕτερός τις --- οὐχ \textbf{ἑόρακα} πώποτε}
Le signe mis pour signifier un changement d'interlocuteur\textsuperscript{1} ayant par la suite été gratté (\latin{ut videtur}).
L'entièreté du vers est effectivement à restituer à Euripide : une scholie écrite juste au-dessus du vers l'atteste : 
\par{\leftskip=1em \textbf{ἐπειδὴ} οὐ πάλαι ἤρξατο διδάσκειν, ἀλλὰ τρισὶν πρὸ τούτων ἔτεσιν.\par 
\textbf{Puisqu}'il n'a pas commencé récemment de faire jouer des pièces, mais depuis trois ans. [sept, en réalité, et selon la façon grecque de compter]}
Cette scholie est évidemment là pour expliquer l'étonnement d'Euripide : οὐχ ἑόρακα\textbf{ς} πώποτε; « Tu ne l'as donc jamais vu ?! [alors que cela fait trois année qu'il monte des pièces] » ; le scholiaste ne l'aurait pas ainsi rédigée (en commençant par ἐπειδὴ...) si, au vers 32, le `οὐχ ἑόρακα πώποτε' était prononcé par Mnésiloque, « je ne l'ai jamais vu ».\par 
Cette interprétation, outre qu'elle est anciennement corroborée par la scholie, est aussi celle qui permet la lecture la plus drôle : οὐχ ἑόρακας πώποτε; --- μῶν ὁ δασυπώγων; --- οὐχ ἑόρακας πώποτε, « Tu ne l'as jamais vu ?! --- N'est-il pas barbu ? --- Tu ne l'as jamais vu !.. », bien meilleure par le rythme et la répétition déconvenue d'Euripide que l'échange premièrement transcrit sur le manuscrit de Ravenne où Mnésiloque commençait par avouer son ignorance : (Μνη.) οὐχ ἑόρακα πώποτε. μῶν ὁ δασυπώγων; --- (Εὐρ.) οὐχ ἑόρακας πώποτε.
\par{\fontsize{7.6}{9.4}\selectfont\leftskip=1em\noindent \textsuperscript{1}\,Non un tiret, \textit{paragraphos}, ni un double point ` : ', dont l'usage est courant dans notre manuscrit lorsque le changement d'interlocuteur intervient au milieu d'un vers, mais en l'occurrence une indication de personnage : μ\auDessus{ν}{Η}.\par}}
\vers[n=31]{}{}{χωρὶς γὰρ αὐτοῖν ἑκατέρου ᾽στὶν ἡ φύσις.}{C'est que distincte est l'essence de chacune de ces choses.}
\vers[n=32]{}{}{χωρὶς γὰρ αὐτοῖν ἑκατέρου ᾽στὶν ἡ φύσις.}{C'est que distincte est l'essence de chacune de ces choses.}
\vers[n=33]{ΕΥ.}{Euripide}{χωρὶς γὰρ αὐτοῖν ἑκατέρου ᾽στὶν ἡ φύσις.}{C'est que distincte est l'essence de chacune de ces choses.}
\vers[n=34]{}{}{χωρὶς γὰρ αὐτοῖν ἑκατέρου ᾽στὶν ἡ φύσις.}{C'est que distincte est l'essence de chacune de ces choses.}
\vers[n=35]{}{}{χωρὶς γὰρ αὐτοῖν ἑκατέρου ᾽στὶν ἡ φύσις.}{C'est que distincte est l'essence de chacune de ces choses.}
\comm{35}{καὶ μὴν βεβίνηκας σύ γ᾽, ἀλλ᾽ οὐκ οἶσθ᾽ ἴσως.}{Le ton est inhabituel -- voire plutôt choquant -- dans la bouche d'Euripide.
Qu'à cela ne tienne : puisqu'il semble d'autre part qu'Agathon tourne le dos lorsqu'il accorde ses faveurs\textsuperscript{1}, nos commentateurs en concluent qu'il est logique que Mnésiloque ne le connaisse pas -- sauf éventuellement au sens biblique... de Sodome, s'entend. 
Et c'est le grave et docte Euripide qui colporterait ce genre de ragots. 
Y a-t-il besoin de commenter ce genre de commentaires ?.. je précise : ces commentaires \textit{modernes} ; car le βεβίνηκας du \textit{Ravennas} n'a pas laissé derrière lui la moindre scholie ancienne ; chose étonnante lorsqu'on sait que les anciens commentateurs étaient presqu'autant portés au scabreux que nos modernes. 
Il faut croire que le texte original d'Aristophane, l'\allemand{Urtext} des \textit{Thesmophories}, comportait une tout autre leçon : une banalité sur laquelle aucun commentateur ancien ne se sera arrêté.
Las !.. on ne détruit que ce que l'on remplace ; à défaut d'avoir une proposition alternative solidement étayée, il faut nous résoudre à \latin{crux}-désespérer...

















































\par{\fontsize{7.6}{9.4}\selectfont\leftskip=1em\noindent \textsuperscript{1}\,Commentaire paraissant de prime abord tout-à-fait fantaisiste et émanant d'une imagination délurée. Mais il semble corroboré par le  \textsuperscript{1}\,Voir le vers 258 et sa scholie ; Agathon se livrerait au travestissement dans l'interlope \textit{Athens by night} pour y être éprouvé à la manière d'une femme -- mais le commentaire déluré du scholiaste n'est peut-être qu'une mauvaise interprétation du présent vers 35 adjointe à une forte dose d'imagination lubrique. Ce qui est sûr, c'est qu'au vers 205, Agathon dit craindre s'il était reconnu d'être lynché pour avoir chapardé nuitamment les œuvres de la Cypris : les Athéniennes voient en lui un concurrent irrégulier.\par}}
\end{texte}
\end{folio}
\end{document}